\section{Context: paco, gpaco, and Interaction Trees}
\label{sec:paco}

\subsection{Improvement 1: paco/gpaco}

\code{paco}~\cite{paco} (\textbf{pa}rametrized \textbf{co}induction) formalizes
Knaster-Tarski\ldots allowing us to find the \code{gfp} in a more permissive
way\ldots by giving us rules for accumulation and composition.

\gap{$\mathrm{G}_f(x) \stackrel{\mathrm{def}}{=} \nu y.\, f(x \sqcup y)$, with
$G$ the ``parameterized'' \code{gfp}, and $\nu f \equiv \mathrm{G}_f(\bot)$.
Then the three rules from slide 39: COMPOSE, for combining theorems;
ACCUMULATE, for adding knowledge during a proof; UNFOLD, for exposing the
monotone function \code{f}.}

\code{gpaco}~\cite{gpaco} (\textbf{g}eneralized \textbf{paco}) exploits \emph{compatible
enhancements} to the \code{gfp} for \textbf{up-to} reasoning.

\textbf{Paco stats:} \gap{The scorecard from slide 41: 1~\checkmark,
2~(partial), 3~\checkmark, 4~\checkmark, 5~\checkmark. \ldots but \underline{much}
better than \code{cofix}!}

\subsection{From paco, Interaction Trees (ITrees)~\cite{itrees}}

Idea: Abstract away \underline{all} coinductive proof, and do reasoning modulo an
equational theory.

ITree: A data structure for:
\begin{enumerate}
  \item Modeling \underline{infinite} and \underline{impure computation}
  \item \underline{Hiding the low-level coinductive reasoning} one would do with
        \code{paco}.
  \item \underline{Extractable} execution for verified software
\end{enumerate}

\code{itree} is the \code{CoInductive itreeF}.

Paco is pervasive in ITrees, and is used to define bisimilarity.
\gap{The search hit from slide 44, ``667 results in 54 files'', beside
\code{eqitF} as the \code{b} and \code{eqit} as the \code{gfp b}. This is the
slide that justifies calling the port large.}

\subsection{\ldots but paco has some weaknesses}

\begin{itemize}
  \item Long compatibility proofs
  \item Lemmas reasoning about compatibility are hard to read and maintain
  \item Lots of individual tactics make it a bit tricky to use and automate
        (6 paco-specific tactics in this proof)
\end{itemize}

\gap{The three code figures from slides 46, 47 and 48. At least one has to
survive into the paper; the 6-tactics-in-one-proof example on slide 48 is the
one that carries a number.}

\subsection{Can we do better?}

Paco is a complete theory of up-to coinduction: Problem is only in usability.

\textbf{Main problems:}
\begin{enumerate}
  \item Compatibility proofs can get long,
  \item Proofs with \code{gpaco} closures can get confusing and use different
        techniques,
  \item Lots of tactics to remember (4-8+ unique paco tactics per proof)
  \item Minor wart: ``arity hell:'' each paco operator is re-defined for
        different arities, e.g., \code{paco1}, \code{paco2}, up to
        \code{paco16}\ldots (also makes portable automation harder to write)
\end{enumerate}

So, our goal is to find a theory that gives ``nicer'' proofs, is easier to
maintain and teach, etc.
